\section{Conclusion}
\label{sec:conclusion}

We examined the dependency structure of Brazilian equities traded on B3 using an econophysics workflow. The raw historical universe covers 1998--2025; the RMT, clustering and network findings are based on its complete 58-asset panel from 23 August 2006 to 15 August 2024. The main findings can be summarized along four dimensions.

First, the stylized-fact and correlation analysis shows that Brazilian equities display signatures commonly associated with complex financial systems: heavy-tailed return behaviour, volatility clustering, persistent absolute-return dependence and a moderate but positive average cross-correlation ($\bar{\rho}\approx0.24$). Within-sector correlations are larger than between-sector correlations on average, confirming that sectoral organization is reflected in the dependency matrix.

Second, Random Matrix Theory decomposes the correlation matrix. Five eigenvalues exceed the Mar\v{c}enko--Pastur upper bound, with the largest eigenvalue accounting for approximately 37.3\% of the trace of the correlation matrix. The leading component captures a market-mode component, while the remaining eigenvectors outside the random band suggest commodity, financial and utility-related group structures. The reconstruction checks indicate numerical accuracy within floating-point precision, and the group-mode matrix retains localized dependency patterns after the dominant market component is removed.

Third, financial network analysis through MST and PMFG representations shows that RMT filtering changes the structure of asset dependencies. The group-mode PMFG has lower mean edge correlation but higher clustering than the original PMFG, indicating stronger local clustering after market-mode removal. Hub rankings shift from networks dominated by broad market co-movement to networks where different assets occupy central positions after excluding the leading component. The PMFG retains 168 edges with triangles and 4-cliques, leaving local graph structures that are absent from the MST and useful for subsector-level analysis.

Fourth, the chronological volatility-modelling experiment identifies a limited ex-post association between network topology and model performance. Models augmented with full-sample network descriptors obtain the lowest average QLIKE in selected model--horizon comparisons, while the CNN-1D model is competitive in MAE, RMSE and held-out $R^2$. Lagged-volatility features remain dominant, and selected PMFG centralities have smaller importance. These comparisons do not establish a real-time forecasting benefit because the RMT and network descriptors use full-sample information.

We view the net contribution of network descriptors as methodological and incremental: they do not replace lagged-volatility dynamics, but they summarize structural information from the cross-sectional geometry of the market. The essential next step is to construct RMT and network features recursively at each forecast origin, then repeat model selection and evaluation without future information. That design, together with a broader target-asset set and richer temporal neural or graph architectures, can determine whether the structural association survives as a real-time forecasting benefit.
